%% 第 5 章：代数基本定理及其在复空间算子上的推论
%% 本文件由 tools/md2tex.py 自动生成，请勿直接编辑。

\chapter{代数基本定理及其在复空间算子上的推论}


\begin{jthm}{代数基本定理}{r:6:1}
任意非常数复系数多项式：

\[
p(z)=a_nz^n+\cdots+a_1z+a_0,
\qquad
a_n\ne0,
\]

至少有一个复根。

等价地，任意 $n$ 次复系数多项式都可以完全分解为一次因子的乘积：

\[
p(z)
=
a_n(z-\lambda_1)\cdots(z-\lambda_n).
\]

其中 $\lambda_1,\ldots,\lambda_n$ 允许重复。重复次数称为相应根的重数。
\end{jthm}

\section{为什么它对算子重要}

设 $V$ 是一个非零有限维复向量空间，且：

\[
T:V\to V.
\]

$T$ 的特征多项式：

\[
p_T(z)=\det(zI-T)
\]

是一个次数为：

\[
\dim V
\]

的复系数多项式。

因此，由代数基本定理，$p_T$ 至少有一个复根 $\lambda$：

\[
p_T(\lambda)=0.
\]

这意味着：

\[
\det(\lambda I-T)=0.
\]

所以：

\[
\lambda I-T
\]

不可逆；等价地：

\[
T-\lambda I
\]

的零空间非平凡。

因此，存在非零向量 $v$ 满足：

\[
(T-\lambda I)v=0.
\]

也就是：

\[
T(v)=\lambda v.
\]

所以得到结论：

\begin{jthm}{复向量空间上的算子一定存在特征值}{r:6:2}
若 $V$ 是非零有限维复向量空间，则任意算子：

\[
T:V\to V
\]

至少存在一个特征值。
\end{jthm}

\section{这个推论的意义}

这保证了复空间上的任意算子至少存在一个一维不变子空间：

\[
\operatorname{span}(v).
\]

因此，我们总能先从一个特征向量出发，再通过限制算子、商算子或归纳过程继续分析其余空间。

这正是复数域上任意算子都可以上三角化的起点。
